%------------------------------------------------------------------------------%
\documentclass[fleqn,utf8,aspectratio=169,12pt]{beamer}

\usepackage{integracao_e_series}

\title{Séries Numéricas -- Teste da Integral}

%------------------------------------------------------------------------------%
\begin{document}

\addtitle

%------------------------------------------------------------------------------%
\section{Série com Termos Não-Negativos}

%------------------------------------------------------------------------------%
\begin{frame}{Recapitulando}
  Sequências numéricas
  \[
    \big( a_k \big)_{k=1}^\infty \qquad\qquad\qquad a_1,\; a_2,\; a_3,\; \dots
  \]
  \pause
  Séries numéricas
  \[
    \sum_{k=1}^\infty a_k = a_1 + a_2 + a_3 + \cdots + a_k + \cdots
  \]
  \pause
  Somas parciais
  \[
    S_n = \sum_{k=1}^n a_k = a_1 + a_2 + a_3 + \cdots + a_n
  \]
\end{frame}

%------------------------------------------------------------------------------%
\begin{frame}{Série com Termos Não-Negativos}
  Caso particular de séries onde $\quad a_n \geq 0 \quad \forall n$

  \pause
  \bigskip
  Uma série
  \[\sum_{n=1}^\infty a_n \qquad \text{com} \qquad a_n\geq 0 \]

  \pause
  \bigskip
  é \emph{convergente} se, e somente se, sua sequência de somas parciais,  \\
  $S_n$, é uma sequência \emph{limitada superiormente}.
\end{frame}

%------------------------------------------------------------------------------%
\section{Teste da Integral}

%------------------------------------------------------------------------------%
\begin{frame}{Teste da Integral}
  Trocamos o problema de verificar se uma série converge para o de
  verificar se uma integral converge
  \[
    \sum_{n=1}^{\infty} a_n  \;\text{ converge?}
    \qquad \text{ por } \qquad
    \int_{1}^{\infty} f(x) \,dx  \;\text{ converge?}
  \]

  \pause
  O valor da integral \emph{\textbf{não é}} uma boa aproximação para a soma da série
  \[ \sum_{n=1}^{\infty} a_n  \quad \not\approx \quad \int_{1}^{\infty} f(x) \,dx\]
\end{frame}

%------------------------------------------------------------------------------%
\begin{frame}{Diferença Entre a Integral e a Série}
  \begin{center}
    % Title: gl2ps_renderer figure
% Creator: GL2PS 1.4.2, (C) 1999-2020 C. Geuzaine
% For: Octave
% CreationDate: Sat Sep 18 17:28:01 2021
\setlength{\unitlength}{1pt}
\begin{picture}(0,0)
\includegraphics[width=14cm,height=5cm]{fig_teste_integral_comparacao}
\end{picture}%
\begin{picture}(396,141)(0,0)
\fontsize{10}{0}\selectfont\put(56.0551,9.94382){\makebox(0,0)[t]{\textcolor[rgb]{0.15,0.15,0.15}{{1}}}}
\fontsize{10}{0}\selectfont\put(125.504,9.94382){\makebox(0,0)[t]{\textcolor[rgb]{0.15,0.15,0.15}{{3}}}}
\fontsize{10}{0}\selectfont\put(194.953,9.94382){\makebox(0,0)[t]{\textcolor[rgb]{0.15,0.15,0.15}{{5}}}}
\fontsize{10}{0}\selectfont\put(264.402,9.94382){\makebox(0,0)[t]{\textcolor[rgb]{0.15,0.15,0.15}{{7}}}}
\fontsize{10}{0}\selectfont\put(333.85,9.94382){\makebox(0,0)[t]{\textcolor[rgb]{0.15,0.15,0.15}{{9}}}}
\fontsize{10}{0}\selectfont\put(16.3318,72.9685){\makebox(0,0)[r]{\textcolor[rgb]{0.15,0.15,0.15}{{0,5}}}}
\fontsize{10}{0}\selectfont\put(16.3318,128.48){\makebox(0,0)[r]{\textcolor[rgb]{0.15,0.15,0.15}{{1,0}}}}
\end{picture}

  \end{center}
\end{frame}

%------------------------------------------------------------------------------%
\begin{frame}[t]{Teste da Integral}
  \onslide<+->{Seja $(a_n)$ uma sequência tal que}
  \begin{enumerate}[<+->]
    \item seus termos são positivos, $a_n>0$
    \item $a_n=f(n)$
    \item \emph{a função $f$} é contínua, positiva e decrescente para todo $x \geq N$
  \end{enumerate}

	\bigskip
  \onslide<+->{Então, a série \tabto{26mm}\(\sum_{n=N}^\infty a_n\) \tabto{50mm} converge}

  \smallskip
	\onslide<+->{\emph{se, e somente se}}

  \medskip
  \onslide<+->{a integral     \tabto{26mm}\(\int_N^\infty f(x)dx \) \tabto{50mm} converge}
\end{frame}

%------------------------------------------------------------------------------%
\begin{frame}{Ilustração do Teste da Integral -- 1}
  \begin{center}
    \scalebox{.9}{% Title: gl2ps_renderer figure
% Creator: GL2PS 1.4.2, (C) 1999-2020 C. Geuzaine
% For: Octave
% CreationDate: Fri Sep 17 21:56:23 2021
\setlength{\unitlength}{1pt}
\begin{picture}(0,0)
\includegraphics[width=16cm,height=5cm]{fig_teste_integral_right}
\end{picture}%
\begin{picture}(453,141)(0,0)
\fontsize{10}{0}\selectfont\put(120.945,10.9531){\makebox(0,0)[t]{\textcolor[rgb]{0.15,0.15,0.15}{{$N\!-\!1$}}}}
\fontsize{10}{0}\selectfont\put(156.22,10.9531){\makebox(0,0)[t]{\textcolor[rgb]{0.15,0.15,0.15}{{$N$}}}}
\fontsize{10}{0}\selectfont\put(191.496,10.9531){\makebox(0,0)[t]{\textcolor[rgb]{0.15,0.15,0.15}{{$N\!+\!1$}}}}
\fontsize{10}{0}\selectfont\put(262.047,10.9531){\makebox(0,0)[t]{\textcolor[rgb]{0.15,0.15,0.15}{{$N\!+\!3$}}}}
\fontsize{10}{0}\selectfont\put(332.598,10.9531){\makebox(0,0)[t]{\textcolor[rgb]{0.15,0.15,0.15}{{$N\!+\!5$}}}}
\fontsize{10}{0}\selectfont\put(403.15,10.9531){\makebox(0,0)[t]{\textcolor[rgb]{0.15,0.15,0.15}{{$N\!+\!7$}}}}
\end{picture}
}
  \end{center}

  \pause
  A série é \emph{``maior''} do que a integral
  \pause
  \(
    \qquad\qquad
    \int_{N}^{\infty} f(x)\,dx
    \quad\leq\quad
    \sum_{n=N}^{\infty} a_n
  \)
\end{frame}

%------------------------------------------------------------------------------%
\begin{frame}{Ilustração do Teste da Integral -- 2}
  \begin{center}
    \scalebox{.9}{\input{fig_teste_integral_left}}
  \end{center}

  \pause
  A série é \emph{``menor''} do que a integral
  \pause
  \(
    \qquad\qquad
    \sum_{n=N}^\infty a_k
    \quad\leq\quad
    a_{N} + \int_{N}^\infty f(x)\,dx
  \)
\end{frame}

%------------------------------------------------------------------------------%
\begin{frame}{Comparação entre a Série e a Integral}
  Temos portanto as relações
  \[
    \int_{N}^{\infty} f(x)\,dx
    \quad\leq\quad
    \sum_{k=N}^{\infty} a_k
    \quad\leq\quad
    a_N + \int_{N}^{\infty} f(x)\,dx
  \]
\end{frame}

%------------------------------------------------------------------------------%
\begin{frame}{Lembrando como Calcular uma Integral Imprópria}
  \[
    \int_{a}^{\infty} f(x)\,dx
    \pause \;=\;
    \lim_{b\to\infty} \; \int_{a}^{b} f(x)\,dx
    \pause \;=\;
    \lim_{b\to\infty} \left( F(x) \evalat{a}{b} \; \right)
    \pause \;=\;
    \lim_{b\to\infty} \big( F(b) - F(a) \big)
  \]
\end{frame}

%------------------------------------------------------------------------------%
\section{Exemplos}

%------------------------------------------------------------------------------%
\begin{question}
  Use o teste da integral para determinar se a série
  \[
  \sum_{n=1}^{\infty} \frac{1}{n^2}
  \]
  converge ou diverge

  \pause
  \bigskip
  Certifique-se de que as condições do teste da integral sejam satisfeitas
\end{question}

%------------------------------------------------------------------------------%
\begin{resolution}{Solução}
  Podemos aplicar o teste da integral, pois, a \emph{função}
  \(
    f(x) = \frac{1}{x^2},
  \)

  \pause
  para $x\geq 1$,
  \pause
  é \emph{positiva},
  \pause
  \emph{contínua}
  \pause
  e \emph{decrescente}
\end{resolution}

%------------------------------------------------------------------------------%
\begin{resolution}{Solução}
  \begin{align*}
    \onslide<+->{\int_{1}^{\infty} \frac{1}{x^2}\,dx}
    \onslide<+->{& = \lim_{b\to\infty} \int_{1}^{b} \frac{1}{x^2}\,dx           \\[2mm]}
    \onslide<+->{& = \lim_{b\to\infty} \left( -\frac{1}{x} \right)\evalat{1}{b} \\[2mm]}
    \onslide<+->{& = \lim_{b\to\infty} \left( -\frac{1}{b} +1 \right)}
    \onslide<+->{ \;=\; 1}
  \end{align*}

  \onslide<+->{Como a integral converge a série também converge}
\end{resolution}

%------------------------------------------------------------------------------%
\begin{resolution}{Diferença Entre a Soma e a Integral}
  Note que
  \[
    \sum_{n=1}^{\infty} \frac{1}{n^2}
    \;=\;       \frac{\pi^2}{6}
    \;\approx\; 1,6449
    \;\neq\;    1
  \]

  \vspace{\baselineskip}
  O cálculo da soma da série está na Seção 9.6 -- Problema de Basileia
\end{resolution}

%------------------------------------------------------------------------------%

%------------------------------------------------------------------------------%
\begin{question}
  Use o teste da integral para determinar se a série converge ou diverge.
  Certifique-se de que as condições do teste da integral sejam satisfeitas.
  \[
    \sum_{n=1}^{\infty} \frac{n}{n^2+4}
  \]
\end{question}

%------------------------------------------------------------------------------%
\begin{resolution}{Solução}
  A função \(f(x)=\frac{x}{x^2+4}\) é \emph{positiva} e \emph{contínua} para $x\geq 1$

  \pause
  \medskip
  Para verificar que ela é \emph{decrescente}, calculamos sua derivada
  \[
    f'(x)
    \pause
    = \frac{1(x^2+4) - x(2x)}{(x^2+4)^2}
    \pause
    = \frac{4-x^2}{(x^2+4)^2}
  \]

  \pause
  $f'(x)<0$, e, portanto, $f$ é decrescente, para todo $x>2$

  \pause
  \medskip
  Precisamos de um $N$ inteiro a partir do qual $f$ seja decrescente

  \pause
  \medskip
  Escolhemos $N=3$
\end{resolution}

%------------------------------------------------------------------------------%
\begin{resolution}{Solução}
  \onslide<+->{Podemos, agora, aplicar o teste da integral}

  \vspace{\baselineskip}
  \onslide<+->{Primeiro calculamos a primitiva de \(f(x)=\frac{x}{x^2+4}\)}
  \onslide<+->{usando substituição
  \[
    u=x^2+4 \qquad \text{e} \qquad du=2xdx
  \]}
  \vspace{-\baselineskip}
  \begin{align*}
    \onslide<+->{F(x) = \int\frac{x}{x^2+4}\,dx
                 & = \int \frac{1}{2u}\,du      \\[2mm]}
    \onslide<+->{& = \frac{1}{2}\ln u + C       \\[2mm]}
    \onslide<+->{& = \frac{1}{2}\ln \left(x^2+4\right) + C}
  \end{align*}
\end{resolution}

%------------------------------------------------------------------------------%
\begin{resolution}{Solução}
  \onslide<+->{Calculando a integral imprópria temos}
  \begin{align*}
    \onslide<+->{\int_{3}^{\infty}
    & \frac{x}{x^2+4}\,dx = \lim_{b\to\infty} \int_{3}^{b} \! \frac{x}{x^2+4}\,dx  \\}
    \onslide<+->{& = \lim_{b\to\infty} \left[ F(x) \right]\evalat{3}{b}       \\}
    \onslide<+->{& = \lim_{b\to\infty} \left[ \frac{1}{2} \ln\left(x^2+4\right) \right]\evalat{3}{b}       \\}
    \onslide<+->{& = \lim_{b\to\infty} \left[ \frac{\; \ln\left(b^2+4\right) -\ln\left(3^2+4\right) \:}{2} \right]}
    \onslide<+->{\;=\; \infty}
  \end{align*}
  \onslide<+->{Como a integral diverge a série também diverge}
\end{resolution}

%------------------------------------------------------------------------------%


%------------------------------------------------------------------------------%
\section{Estimativa de Erro}

%------------------------------------------------------------------------------%
\begin{frame}{Estimativa de Erro}
  Não temos uma fórmula para a soma da série

  \pause
  \vspace{\baselineskip}
  Podemos usar as somas parciais para aproximar a soma

  \pause
  \vspace{\baselineskip}
  Quantos termos precisamos?
\end{frame}

%------------------------------------------------------------------------------%
\begin{frame}{Resto ou Erro}
  \begin{align*}
    \action<+->{S & = \sum_{k=1}^{\infty} a_k} \\[2mm]
    \action<+->{  & = \sum_{k=1}^{n} a_k + \sum_{k=n+1}^{\infty} a_k} \\[2mm]
    \action<+->{  & = S_n + R_n}
  \end{align*}
  \action<+->{
  $R_n$ é o \emph{resto da série}, isso é, o que faltou ser calculado
  \[ R_n = \sum_{k=n+1}^{\infty} a_k \]}
\end{frame}

%------------------------------------------------------------------------------%
\begin{frame}{Estimativa de Erro}
  Se uma série é \emph{convergente pelo teste da integral} podemos estimar o resto
  \[ R_n = a_{n+1} + a_{n+2} + a_{n+3} + a_{n+4} + \cdots \]
  \pause
  comparando a área dos retângulos com a da integral temos
  \pause
  \[ R_n \geq \int_{n+1}^{\infty} f(x)\,dx \]
  \pause
  \[ R_n \leq \int_{n}^{\infty} f(x)\,dx \]
\end{frame}

%------------------------------------------------------------------------------%
\begin{frame}{Limitantes para o resto no teste da integral}
  Se a série
  \(
    \quad \sum_{k=1}^{\infty} a_k
  \)

  \vspace{0.5\baselineskip}
  converge segundo o teste da integral, então o resto satisfaz
  \[
    \int_{n+1}^{\infty} f(x)\,dx
    \quad\leq\quad
    R_n
    \quad\leq\quad
    \int_{n}^{\infty} f(x)\,dx
  \]
\end{frame}

%------------------------------------------------------------------------------%
\section{Exemplos}

%------------------------------------------------------------------------------%
\begin{question}
  Sabemos que a série
  \[
    \sum_{n=1}^\infty\frac{1}{\,n^2\,}
  \]
  converge, pois é uma série $p$, com $p=2$

  \vspace{\baselineskip}
  Descubra $n$ tal que o resto $R_n$ da aproximação $S_n$ seja menor que \num{0.01}
\end{question}

%------------------------------------------------------------------------------%
\begin{resolution}{Solução}
  Sabemos que o Teste da Integral pode ser aplicado, pois
  \[
    \int_n^\infty\frac{\,dx\,}{x^2}
    = \lim_{b\to\infty}\left[ -\frac{1}{\,b\,} + \frac{1}{\,n\,} \right]
    = \frac{1}{\,n\,}
  \]

  \pause
  \vspace{\baselineskip}
  Assim
  \[
    \frac{1}{\,n+1\,} \;\leq\; R_n \;\leq\; \frac{1}{\,n\,}
  \]
\end{resolution}

%------------------------------------------------------------------------------%
\begin{resolution}{Solução}
  Para garantir que $R_n<\num{0.01}$ impomos que
  \[
    \frac{1}{\,n\,} < \num{0.01}
  \]
  isolando $n$ obtemos $n>100$

  \pause
  \vspace{\baselineskip}
  A aproximação $S_{100}$ tem um erro menor do que \num{0.01}
\end{resolution}

%------------------------------------------------------------------------------%

%------------------------------------------------------------------------------%
\section{Lista Mínima}

%------------------------------------------------------------------------------%
\begin{frame}{Lista Mínima}
  Estudar as Seção 6.4 da Apostila

  \vspace{\baselineskip}
  Exercícios: 3, 7a, 8, 9, 10

  \vfill
  Atenção: A prova é baseada no livro, não nas apresentações
\end{frame}

%------------------------------------------------------------------------------%
\end{document}
%------------------------------------------------------------------------------%
